\section{总结与延伸}

\subsection{核心结论}

\begin{enumerate}
\item 高继扬的方法论是归纳总结、目标倒推和概率管理，这种方法从考试、论文、工程到创业一以贯之。
\item Waymo 训练系统工程：拆解、测量、infra、长期执行；但离客户和经营较远。
\item Momenta 训练量产交付：从 demo 到产品，从 research lab 到客户价值，组织必须被真实项目洗礼。
\item 自动驾驶商业模式的关键是数据和客户：谁拥有车辆、用户和利润池，谁就更容易形成数据闭环。
\item 星海图选择整机，是因为整机是数据和客户价值的入口，而不是因为硬件本身更浪漫。
\item Data Recipe 是机器人公司把真实世界数据转成模型能力的核心机制。
\item 机器人大脑不是单个 LLM，而是 VLM 上层拆解与 VLA 动作模型的双系统，并受端侧延迟约束。
\item 算法创新重要，但在开源时代传播周期短；长期壁垒更可能来自整机、供应链、数据、客户和组织。
\item “到土里去”说明机器人创业的理想主义必须以务实、ROI、交付和客户现场为形态。
\end{enumerate}

\subsection{开放问题}

最后保留开放问题，是因为具身智能仍处在路线未收敛的阶段。本期给出了强烈的现实主义框架，但框架如何在未来几年的产品和模型中兑现，仍要继续观察。

\begin{itemize}
\item 具身智能公司到底应该先做开发者市场，还是直接进入生产力场景？
\item VLA 的通用性会来自模型规模、数据规模，还是来自整机和场景闭环？
\item 有整机能力的创业公司，能否在数据上形成相对大厂的长期优势？
\item 当算法传播周期越来越短，机器人公司的护城河会不会越来越向供应链、客户和数据 recipe 转移？
\item 中国供应链优势能否系统性转化为具身智能基础模型优势？
\end{itemize}

\subsection{拓展阅读}

\begin{itemize}
\item EP121 对 DeepMind 谭捷的访谈：机器人、跨本体、世界模型和 Gemini Robotics。
\item EP109 与谢晨聊机器人数据荒：仿真、合成数据和数据产业。
\item EP098 机器人基座模型和 VLA 经典论文讲解：理解 VLA 技术脉络。
\item 自动驾驶系统架构、端到端自动驾驶、robotaxi 商业模式相关材料。
\item 机器人供应链、端侧部署、遥操作数据和后训练工具相关资料。
\end{itemize}

\begin{importantbox}{最后的判断}
这期最值得带走的判断是：具身智能不是从“模型很强”自然推出“机器人可用”，而是从整机、客户、数据、模型和组织共同形成闭环。真正的技术浪漫，可能不是站在云端讲未来，而是愿意把头伸进土里，把未来一点点做出来。
\end{importantbox}

\end{document}
